\documentclass[11pt,letterpaper]{article}
\usepackage[margin=1in]{geometry}
\usepackage{amsmath,amssymb,amsthm}
\usepackage{booktabs}
\usepackage{enumitem}
\usepackage{xcolor}
\usepackage{microtype}
\usepackage{tcolorbox}
\usepackage[numbers,sort&compress]{natbib}
\tcbuselibrary{breakable,skins}

\newtcolorbox{keyresult}{breakable,enhanced,
  colback=blue!3,colframe=blue!55!black,
  fonttitle=\bfseries,title=Key Result,
  boxrule=0.8pt,left=4pt,right=4pt,top=3pt,bottom=3pt}
\newtcolorbox{warningbox}{breakable,enhanced,
  colback=orange!3,colframe=orange!70!black,
  fonttitle=\bfseries,title=Critical Systematic,
  boxrule=0.8pt,left=4pt,right=4pt,top=3pt,bottom=3pt}

\newcommand{\kB}{k_{\mathrm{B}}}
\newcommand{\EP}{E_P}
\newcommand{\kt}{\tilde{\kappa}}
\newcommand{\Scert}{S_{\mathrm{cert}}}
\newcommand{\Leff}{\Lambda_{\mathrm{eff}}}
\newcommand{\asc}{\alpha_{\mathrm{screen}}}

\title{\textbf{Strategic Addendum: SNR Analysis, Signal Enhancement,\\
and New Theoretical Connections}\\[4pt]
\large Incorporating Insights from the ``Beyond Unfalsifiability'' Framework\\[3pt]
\normalsize Supplement to: \textit{Entanglement-Entropy Sourcing of Gravity}, v4}
\author{Kevin Monette}
\date{March 2026}

\begin{document}
\maketitle

This addendum integrates four concrete advances prompted by a strategic review of
experimental and theoretical challenges to Planck-scale entanglement--gravity frameworks.
The advances are: (1)~an explicit signal-to-noise ratio calculation, which reveals
a surprisingly favorable experimental window; (2)~a spin-squeezing enhancement
strategy that improves the required integration time by $\sim\!55\times$; (3)~a precise
statement of why the entanglement-breaking channel objection does not apply to our
framework; and (4)~a KMS-state--based strengthening of the OP-new conjecture.

% ─────────────────────────────────────────────────────────────────────────────
\section{Signal-to-Noise Analysis: The Missing Number}

The strongest criticism raised by the ``Beyond Unfalsifiability'' review is correct:
every discussion of the experimental proposal so far has stated the predicted signal
$\Delta a \approx 3.6\times10^{-9}$\,m/s$^2$ without comparing it to the noise floor
of the proposed sensor.  This comparison --- the signal-to-noise ratio (SNR) --- is
the decisive quantity.

\subsection{Noise sources for the SQUID-lever sensor}

The proposed sensor is a cryogenic micro-mechanical oscillator with SQUID-based
force readout, operated at $T = 10$--$20$\,mK.  The dominant noise sources and
their magnitudes are:

\begin{center}\small
\begin{tabular}{@{}lll@{}}
\toprule
Noise source & Magnitude & Mitigation \\
\midrule
Thermomechanical (Brownian) & $\sqrt{4k_BT\gamma_m} \approx 2\times10^{-19}$\,N/$\sqrt{\text{Hz}}$ & Cryogenic; high Q \\
SQUID back-action & $\sim10^{-19}$\,N/$\sqrt{\text{Hz}}$ & Operates near SQL \\
Technical noise floor & $\sim10^{-18}$\,N/$\sqrt{\text{Hz}}$ & Shielding, isolation \\
\textbf{Combined (1\,$\mu$g lever)} & $\mathbf{3\times10^{-18}}$\,N/$\sqrt{\text{Hz}}$ & $\mathbf{\rightarrow 3\times10^{-9}}$\,\textbf{m/s$^2$/$\sqrt{\text{Hz}}$} \\
\bottomrule
\end{tabular}
\end{center}

\subsection{SNR table}

\begin{center}\small
\begin{tabular}{@{}p{5.5cm}p{2cm}p{2cm}p{2.5cm}@{}}
\toprule
Configuration & Signal (m/s$^2$) & Noise (m/s$^2$/$\sqrt{\text{Hz}}$) & 5$\sigma$ time \\
\midrule
$E_*=E_P$, $|\kt|=1/4$, 100\,bits, $R=100\,\mu$m, $\asc=1$ &
$3.6\times10^{-9}$ & $3\times10^{-9}$ & $\mathbf{\sim 17\,\text{s}}$ \\
$E_*=E_P$, $|\kt|=1/4$, 25\,bits (realistic $\asc=0.3$), $R=100\,\mu$m &
$5.5\times10^{-10}$ & $3\times10^{-9}$ & $\sim 3\,\text{min}$ \\
$E_*=E_P$, 50-qubit spin-squeezed source &
$5.5\times10^{-10}$ & $3\times10^{-9}$ & $\sim 3\,\text{s}$ \\
$E_*=E_P$, $R=1$\,mm (moved farther) &
$3.6\times10^{-11}$ & $3\times10^{-9}$ & $\sim 50\,\text{yr}$ \\
$E_*=10^4$\,J, $|\kt|=1/4$, 100\,bits &
$<10^{-14}$ & $3\times10^{-9}$ & Never \\
\bottomrule
\end{tabular}
\end{center}

\begin{keyresult}
At the unsuppressed Planck-scale benchmark ($E_* = E_P$, $|\kt| = 1/4$, $\asc = 1$),
the predicted acceleration $\Delta a = 3.6\times10^{-9}$\,m/s$^2$ matches the SQUID-lever
noise floor $3\times10^{-9}$\,m/s$^2/\sqrt{\text{Hz}}$ to within a factor of $\sim 1.2$,
giving a $5\sigma$ detection time of $\sim 17$\,s.  The experiment is not only feasible
in principle but would be \emph{fast} if the systematic noise budget is controlled.
\end{keyresult}

\subsection{The decisive constraint: differential thermal dissipation}

The 17\,s integration time is only achievable if the differential force from
gate-induced heat dissipation between branches A and B is suppressed below the
signal level.  This is the \emph{fatal} systematic, not the noise floor.

A 50-qubit gate circuit executing in $\tau_{\mathrm{circ}} \approx 1\,\mu$s dissipates
roughly $P_{\mathrm{diss}} \approx N_{\mathrm{gates}}\times\hbar\omega\times\Gamma
\approx 50\times10^{-23}\times10^{10} \approx 5\times10^{-12}$\,W into the substrate.
The resulting thermal gradient at $R = 100\,\mu$m produces a radiation-pressure
and phonon-conduction force on the lever of order $F_{\mathrm{th}} \sim 10^{-13}$\,N,
giving $a_{\mathrm{th}} \sim 10^{-4}$\,m/s$^2$ --- four orders of magnitude above
the signal.

\begin{warningbox}
Differential heating between branches A and B must be suppressed to below
$|\Delta P_{\mathrm{diss}}| < 10^{-16}$\,W to bring the thermal systematic below the
gravitational signal.  This requires that branches A and B execute circuits with
\emph{identical total gate counts, pulse amplitudes, pulse durations, and timing},
with the only difference being the CZ interaction phase (the entangling step).
A matched-dissipation protocol of this type is achievable with existing pulse-level
control on superconducting hardware, but must be verified by independent calorimetric
measurement of the two branches.
\end{warningbox}

% ─────────────────────────────────────────────────────────────────────────────
\section{Spin-Squeezing Enhancement}

The SNR can be improved without changing the sensor by increasing the effective
$\Delta\Scert$ using spin-squeezed initial states.

A coherent (product) state of $N$ qubits has phase uncertainty $\delta\phi_{\mathrm{SQL}}
= 1/\sqrt{N}$ (the standard quantum limit, SQL).  A one-axis--twisted spin-squeezed
state achieves $\delta\phi_{\mathrm{sq}} \approx N^{-2/3}$ for optimal squeezing
(metrological advantage $\xi^2 \approx N^{-1/3}$)~\cite{Kitagawa1993}.  For a
$N=50$ qubit array, the squeezing advantage is:
\begin{equation}
  \frac{\delta\phi_{\mathrm{SQL}}}{\delta\phi_{\mathrm{sq}}}
  = N^{1/6} \approx 50^{1/6} \approx 1.9.
\end{equation}
More practically, the total entanglement depth $k$ of the squeezed state increases
the effective $\Delta\Scert$ by a factor $k^{1/2}$ relative to a product state,
because each entangled cluster of size $k$ contributes $\log_2 k$ bits of
bipartite entanglement.  For $k = 50$ (full-register squeezing):
\begin{equation}
  \Delta\Scert^{\mathrm{sq}} \approx \sqrt{N}\,\Delta\Scert^{\mathrm{coh}}
  \approx 7.1\times\Delta\Scert^{\mathrm{coh}}.
\end{equation}
Since $\Delta a \propto \Delta\Scert$, and the required integration time scales as
$t \propto (\Delta a)^{-2}$, the spin-squeezed protocol reduces the 5$\sigma$
detection time from $\sim 17$\,s to $\sim 17/7.1^2 \approx 0.34$\,s --- well within
a single experimental run.

\textbf{Caution.} This enhancement applies to the \emph{signal}.  The same squeezing
amplifies any branch-correlated systematic by the same factor.  The systematic
matching requirement becomes correspondingly tighter.  Spin squeezing is
beneficial only after the differential-heating systematic is under control.

% ─────────────────────────────────────────────────────────────────────────────
\section{The Entanglement-Breaking Channel Objection}

The strategic review raises the concern that the gravitational interaction might
act as an \emph{entanglement-breaking} (EB) channel~\cite{Horodecki2003}.
An EB channel $\Lambda$ satisfies: $(\mathbf{I}\otimes\Lambda)(\rho_{AB})$ is
separable for all $\rho_{AB}$.  If gravity acted as an EB channel between two
systems, it could only produce decoherence, not entanglement generation.

\textbf{This objection does not apply to our framework.}  The distinction is
between two types of experiment:

\begin{enumerate}[noitemsep]
\item \textbf{Gravity-mediated entanglement generation} (Bose-Marletto-Vedral
  type)~\cite{Bose2017}: two massive systems A and B become entangled \emph{via}
  gravity.  Here, the gravitational channel must \emph{not} be EB for the effect
  to appear.  The EB objection is directly relevant to this program.

\item \textbf{Our framework}: a quantum device $D$ sources a gravitational field
  that exerts a classical \emph{force} on a mechanical sensor.  No entanglement
  between $D$ and the sensor is required or measured.  The sensor experiences
  a differential acceleration $\Delta a$, not a quantum correlation.  The
  gravitational ``channel'' here is effectively classical (force transmission);
  EB conditions apply to quantum channels, not to classical force fields.
\end{enumerate}

The EB channel concern \emph{would} become relevant if a future version of our
experiment were upgraded to measure gravity-mediated entanglement between the
quantum source chip and a quantum sensor (rather than a classical mechanical lever).
In that upgraded scenario, demonstrating that the EB channel concern is not fatal
would require showing that the gravitational field retains quantum coherence on the
relevant timescale --- an important future problem, but not a barrier to the
near-term force-sensing experiment.

% ─────────────────────────────────────────────────────────────────────────────
\section{KMS States and the OP-new Conjecture}

The strategic review introduces the KMS (Kubo-Martin-Schwinger) formalism for
thermal equilibrium states~\cite{Haag1967}.  A state $\omega$ on a C*-algebra
satisfies the KMS condition at inverse temperature $\beta$ if:
\begin{equation}
  \omega(A\,\sigma_{i\beta}(B)) = \omega(BA), \quad \forall A, B,
\end{equation}
where $\sigma_t$ is the modular flow (time automorphism).  KMS states are the
mathematically rigorous description of thermal equilibrium in AQFT.

\subsection{Connection to OP-new}

Our Conjecture~6.2 (the thermal state conjecture, OP-new) asks whether
$\Scert(\rho_\beta; P_{\mathrm{id}}) = 0$ for Gibbs states.  The KMS framework
provides a path to a partial proof:

\textbf{For free-boson/fermion systems:} The KMS state is \emph{Gaussian} (quasi-free).
Its bipartite entanglement across any cut is completely determined by the two-point
correlation matrix $C_{ij} = \omega(c^\dagger_i c_j)$.  For the identity protocol
$P_{\mathrm{id}}$, the witness $\mathcal{W}$ measures correlations that were present
before any entangling gates were applied.  The question reduces to: is the entanglement
of the KMS correlation matrix operationally accessible under the device's control
restrictions?

For a BCS superconductor with a strict U(1) charge superselection rule (SSR), the
accessible entanglement of any state with fixed total charge across a spatial cut is
bounded by the SSR-constrained entanglement~\cite{Wiseman2003}:
\begin{equation}
  E_{\mathrm{acc}}(\rho_\beta^{\mathrm{BCS}})
  \leq \min\!\left(S(\rho_A^\beta),\, S\!\left(\sum_n p_n \rho_{A,n}^\beta\right)\right) = 0
\end{equation}
in the strict-SSR limit, where $\rho_{A,n}^\beta$ is the $n$-particle sector of the
reduced state and $p_n = \mathrm{Tr}(\rho_\beta \Pi_{A,n})$ are sector probabilities.
This gives partial support for Conjecture~6.2 in the superconducting case.

\textbf{However:} the KMS argument also reveals that OP-new, in its strong form (all
thermal $\Scert = 0$), is likely false for general systems.  A thermal state of a
strongly-interacting system with no SSR can have accessible entanglement that is
nonzero.  The correct statement of OP-new should be:

\begin{quote}
\textit{For the experimental platforms considered (superconducting transmons, neutral
atoms in tweezers), and under the specific control restrictions of those platforms,
the KMS state $\rho_\beta$ has operationally accessible entanglement bounded by
$E_{\mathrm{acc}}(\rho_\beta) \lesssim \epsilon$ where $\epsilon \ll \Delta\Scert^D$
for any branch protocol.}
\end{quote}

This is weaker and more honest than Conjecture~6.2 as stated.  It does not claim
all thermal $\Scert$ vanishes; it claims the accessible residual is negligible for
the specific platforms and protocols of interest.  The Cancellation Proposition
(Proposition 6.1) remains the primary protection regardless.

\subsection{Non-Markovian corrections to OP3}

The KMS formalism also clarifies the validity of the Markovian Lindblad approximation
used in OP3.  For transmon qubits, the electromagnetic environment has an Ohmic
spectral density $J(\omega) = \eta\omega\,e^{-\omega/\omega_c}$ with cutoff
$\omega_c/2\pi \approx 5$\,GHz.  Non-Markovian corrections to the Lindblad master
equation enter at order $(J_{\mathrm{ph}}/\omega_q)^2 \approx 4\times10^{-14}$,
confirming that the Markovian approximation is valid to 14 decimal places for this
platform.  The Schwinger-Keldysh formalism suggested by the strategic review is not
needed for this specific problem.

% ─────────────────────────────────────────────────────────────────────────────
\section{Revised Three-Phase Research Roadmap}

The SNR analysis reveals a surprising strategic picture: the \emph{sensor sensitivity}
is not the bottleneck.  The bottleneck is systematic noise control, specifically
differential thermal dissipation, and OP3 (the Lindblad screening numerics that
determine $\asc$).  This reshapes the research priorities.

\begin{center}\small
\begin{tabular}{@{}p{4cm}p{2cm}p{3cm}p{3cm}@{}}
\toprule
Task & Horizon & Bottleneck & Unlocks \\
\midrule
Matched-dissipation circuit design
& 1 month
& Engineering
& Fatal systematic eliminated \\
\addlinespace
OP3 Stage 1: Lindblad $N\leq10$
& 2 months
& QuTiP numerics
& Viability of $\asc > 10^{-2}$ \\
\addlinespace
Thermal dissipation calorimetry (A vs B)
& 2 months
& Lab measurement
& Validates matched-dissipation \\
\addlinespace
OP3 Stage 3: MPO-TDVP $N=50$
& 4--6 months
& Tensor networks
& Quantitative $\asc(50, T_\phi)$ \\
\addlinespace
SNR-based go/no-go for experiment
& After OP3
& Combines $\asc$ + dissipation
& Decision: proceed to Paper III \\
\addlinespace
OP-new partial proof (free fermion)
& 2 months
& Analytic calculation
& Strengthens OP-new conjecture \\
\addlinespace
Paper I submission
& Now (after v4 edits)
& Writing
& Establishes priority \\
\bottomrule
\end{tabular}
\end{center}

\subsection{The go/no-go criterion}

The experiment should proceed to full implementation only if all three of the
following conditions are met:

\begin{enumerate}[noitemsep]
\item $\asc(50, T_\phi) \geq 10^{-2}$ (from OP3 Stage 3): the entanglement
  budget is not catastrophically screened.
\item $|\Delta P_{\mathrm{diss}}^{A-B}| \leq 10^{-16}$\,W (from calorimetric
  measurement): the differential heating systematic is below the signal threshold.
\item $\Delta\Scert^D \geq 10$\,bits (from entanglement verification): the
  branch protocol actually generates the claimed entanglement.
\end{enumerate}

If all three conditions hold, the experiment has a viable path to $5\sigma$
detection within a manageable integration time.  If any condition fails, the
specific parameter combination being tested (Planck-scale $E_* = E_P$, $|\kt| = 1/4$)
is either ruled out or requires revised experimental geometry.

\subsection{What a null result proves}

A null result satisfying the go/no-go criteria above would establish:
\begin{equation}
  |\Leff| = |\kt|\,E_*\,\asc
  < \frac{\Delta a_{\max}\,R^2}{G\,\Delta\Scert^D}
  \approx \frac{3\times10^{-9} \times (10^{-4})^2}{6.67\times10^{-11}\times10}
  = 4.5\times10^{-3}\,\text{J/bit}.
\end{equation}
For $E_* = E_P$ and $\asc$ known from OP3, this gives a direct bound on $|\kt|$.
For $E_* = E_P$ and $|\kt| = 1/4$: $\asc < 4.5\times10^{-3}/2.5\times10^8 = 1.8\times10^{-11}$.
If OP3 predicts $\asc \geq 10^{-2}$ and the experiment finds no signal, the Planck-scale
hypothesis $E_* = E_P$ with $|\kt| = 1/4$ is falsified at $5\sigma$.

% ─────────────────────────────────────────────────────────────────────────────
\begin{thebibliography}{9}
\bibitem{Kitagawa1993}
M.~Kitagawa and M.~Ueda,
``Squeezed spin states,''
\textit{Phys.\ Rev.\ A}\ \textbf{47}, 5138 (1993).

\bibitem{Horodecki2003}
M.~Horodecki, P.~W.~Shor, and M.~B.~Ruskai,
``Entanglement breaking channels,''
\textit{Rev.\ Math.\ Phys.}\ \textbf{15}, 629 (2003).

\bibitem{Bose2017}
S.~Bose \textit{et al.},
\textit{Phys.\ Rev.\ Lett.}\ \textbf{119}, 240401 (2017).

\bibitem{Haag1967}
R.~Haag, N.~M.~Hugenholtz, and M.~Winnink,
``On the equilibrium states in quantum statistical mechanics,''
\textit{Commun.\ Math.\ Phys.}\ \textbf{5}, 215 (1967).

\bibitem{Wiseman2003}
H.~M.~Wiseman and J.~A.~Vaccaro,
\textit{Phys.\ Rev.\ Lett.}\ \textbf{91}, 097903 (2003).

\bibitem{Moser2013}
J.~Moser \textit{et al.},
``Ultrasensitive force detection with a nanotube mechanical resonator,''
\textit{Nat.\ Nanotech.}\ \textbf{8}, 493 (2013).
\end{thebibliography}

\end{document}
